\paragraph{全微态阈值的极端矩恢复、冻结压力分解与零温地基层分裂}
在全微态精确反演的有限层比特阈值、冻结相的双指数分离、原子 Euler 因子的 primitive 剥离以及零温四态地基层闭式已经建立之后，还可进一步抽取出下列六条跨模块后果。

\begin{corollary}[全可逆阈值的极端矩重建]\label{cor:xi-full-inversion-threshold-extreme-moment-recovery}
沿用定理 \ref{thm:xi-full-microstate-exact-inversion-bitrate-threshold} 的记号。对每个固定分辨率 \(m\) 与每个 \(q>0\)，记
\[
S_q(m):=\sum_{x\in X_m}d_m(x)^q.
\]
则有
\[
\boxed{\;
M_m=\lim_{q\to\infty}S_q(m)^{1/q}.\;}
\]
从而
\[
\boxed{\;
B_m^{\mathrm{full}}
=
\left\lceil
\log_2\!\left(\lim_{q\to\infty}S_q(m)^{1/q}\right)
\right\rceil.\;}
\]
因此，全微态精确反演阈值不仅由最大纤维规模 \(M_m\) 唯一确定，也由正阶幂和矩谱 \(\{S_q(m)\}_{q>0}\) 的端点极限唯一恢复。
\end{corollary}
\begin{proof}
对固定 \(m\)，向量
\[
\mathbf d^{(m)}:=\bigl(d_m(x)\bigr)_{x\in X_m}
\]
只有有限个坐标。有限维范数极限定理给出
\[
\lim_{q\to\infty}\left(\sum_{x\in X_m}d_m(x)^q\right)^{1/q}
=
\max_{x\in X_m}d_m(x)
=
M_m.
\]
这便得到第一式。再将其代入定理 \ref{thm:xi-full-microstate-exact-inversion-bitrate-threshold} 的恒等式
\[
B_m^{\mathrm{full}}=\left\lceil \log_2 M_m\right\rceil
\]
即可得到第二式。证毕。
\end{proof}

\begin{theorem}[冻结相压力的地址成本--残余简并分解]\label{thm:xi-fixed-freezing-pressure-address-cost-residual-decomposition}
沿用定理 \ref{thm:xi-full-microstate-exact-inversion-bitrate-threshold}、定理 \ref{thm:fold-pressure-freezing-threshold} 与定理 \ref{thm:xi-fixed-freezing-escort-renyi-spectrum-collapse} 的记号。对任意 \(a>a_{\mathrm c}\)，记
\[
\beta_{\mathrm{inv}}
:=
\lim_{m\to\infty}\frac{B_m^{\mathrm{full}}}{m}
=
\frac{\alpha_*}{\log 2},
\]
并记冻结相的 min-entropy 率为
\[
h_\infty(a)
:=
\lim_{m\to\infty}\frac{1}{m}
\Bigl(-\log\max_{x\in X_m}\pi_{m,a}(x)\Bigr)
=
g_*.
\]
则冻结相压力满足严格分解
\[
\boxed{\;
P_a
=
a(\log 2)\,\beta_{\mathrm{inv}}+h_\infty(a).\;}
\]
换言之，冻结相压力恰分解为“实现全微态精确反演所需的指数地址成本”与“冻结端面的残余简并熵率”两部分。
\end{theorem}
\begin{proof}
由定理 \ref{thm:fold-pressure-freezing-threshold}，对任意 \(a>a_{\mathrm c}\) 都有
\[
P_a=a\alpha_*+g_*.
\]
另一方面，定理 \ref{thm:xi-full-microstate-exact-inversion-bitrate-threshold} 给出
\[
\beta_{\mathrm{inv}}=\frac{\alpha_*}{\log 2},
\]
而定理 \ref{thm:xi-fixed-freezing-escort-renyi-spectrum-collapse} 在 \(q=\infty\) 时给出
\[
h_\infty(a)=g_*.
\]
将这两式代回即可得到
\[
P_a
=
a\alpha_*+g_*
=
a(\log 2)\,\beta_{\mathrm{inv}}+h_\infty(a).
\]
证毕。
\end{proof}
\leanverified{paper\_xi\_fixed\_freezing\_pressure\_address\_cost\_residual\_decomposition}

\begin{corollary}[冻结相的斜率--截距反演公式]\label{cor:xi-fixed-freezing-slope-intercept-inversion}
沿用定理 \ref{thm:xi-fixed-freezing-pressure-address-cost-residual-decomposition} 的记号。则有
\[
\boxed{\;
\lim_{q\to\infty}\frac{P_q}{q}=\alpha_*,
\qquad
\beta_{\mathrm{inv}}
=
\frac{1}{\log 2}\lim_{q\to\infty}\frac{P_q}{q}.\;}
\]
并且对任意 \(a>a_{\mathrm c}\)，成立
\[
\boxed{\;
h_\infty(a)
=
P_a-a\lim_{q\to\infty}\frac{P_q}{q}.\;}
\]
因此，冻结相的截距并非独立参数，而是由压力谱减去其无穷远斜率后精确恢复。
\end{corollary}
\begin{proof}
由定理 \ref{thm:xi-time-part57b-pressure-spectrum-slope-intercept-recovery}，
\[
\lim_{q\to\infty}\frac{P_q}{q}=\alpha_*.
\]
再与
\[
\beta_{\mathrm{inv}}=\frac{\alpha_*}{\log 2}
\]
合并，便得第二式。

另一方面，定理
\ref{thm:xi-fixed-freezing-pressure-address-cost-residual-decomposition}
给出
\[
P_a=a\alpha_*+h_\infty(a).
\]
将 \(\alpha_*\) 以
\[
\lim_{q\to\infty}\frac{P_q}{q}
\]
替换，即得
\[
h_\infty(a)
=
P_a-a\lim_{q\to\infty}\frac{P_q}{q}.
\]
证毕。
\end{proof}

\begin{theorem}[冻结相中高阶幂和比值的指数锁定]\label{thm:xi-fixed-freezing-power-sum-ratio-locking}
在严格基态间隙
\[
\alpha_*^{(2)}<\alpha_*
\]
下，沿用定理 \ref{thm:xi-fixed-freezing-escort-renyi-spectrum-collapse} 的记号。记
\[
M_m:=\max_{x\in X_m}d_m(x),
\qquad
K_m:=\#\{x\in X_m:\ d_m(x)=M_m\},
\]
\[
M_m^{(2)}:=\max\{d_m(x):\ d_m(x)<M_m\},
\qquad
S_t(m):=\sum_{x\in X_m}d_m(x)^t.
\]
若次大纤维不存在，则约定 \(M_m^{(2)}:=1\)。对任意固定整数 \(k\ge 1\) 与任意
\[
a>a_{\mathrm c}:=\frac{\log\varphi-g_*}{\alpha_*-\alpha_*^{(2)}},
\]
都有
\[
\boxed{\;
\frac{S_{a+k}(m)}{S_a(m)}
=
M_m^k\Bigl(1+O\bigl(e^{-m\Delta(a)+o(m)}\bigr)\Bigr),\;}
\]
其中
\[
\Delta(a):=a(\alpha_*-\alpha_*^{(2)})-(\log\varphi-g_*)>0.
\]
\end{theorem}
\begin{proof}
把
\[
R_t(m):=\sum_{x:\,d_m(x)<M_m}d_m(x)^t
\]
记作基态外余项，则对任意 \(t>0\) 都有精确分解
\[
S_t(m)=K_mM_m^t+R_t(m).
\]
并且
\[
0\le R_t(m)\le |X_m|\,(M_m^{(2)})^t.
\]
由
\[
M_m=e^{\alpha_* m+o(m)},
\qquad
K_m=e^{g_* m+o(m)},
\qquad
|X_m|=e^{(\log\varphi)m+o(m)},
\]
以及
\[
\limsup_{m\to\infty}\frac1m\log M_m^{(2)}=\alpha_*^{(2)},
\]
可得
\[
\frac{R_a(m)}{K_mM_m^a}
\le
\exp\!\bigl(-m\Delta(a)+o(m)\bigr).
\]
因此
\[
S_a(m)=K_mM_m^a\Bigl(1+O\bigl(e^{-m\Delta(a)+o(m)}\bigr)\Bigr).
\]
同理，
\[
\frac{R_{a+k}(m)}{K_mM_m^{a+k}}
\le
\exp\!\bigl(-m\Delta(a+k)+o(m)\bigr)
\le
\exp\!\bigl(-m\Delta(a)+o(m)\bigr),
\]
因为
\[
\Delta(a+k)=\Delta(a)+k(\alpha_*-\alpha_*^{(2)})\ge \Delta(a).
\]
故
\[
S_{a+k}(m)=K_mM_m^{a+k}\Bigl(1+O\bigl(e^{-m\Delta(a)+o(m)}\bigr)\Bigr).
\]
两式相除便得到
\[
\frac{S_{a+k}(m)}{S_a(m)}
=
M_m^k\,
\frac{1+O\bigl(e^{-m\Delta(a)+o(m)}\bigr)}
{1+O\bigl(e^{-m\Delta(a)+o(m)}\bigr)}
=
M_m^k\Bigl(1+O\bigl(e^{-m\Delta(a)+o(m)}\bigr)\Bigr).
\]
证毕。
\end{proof}

\begin{corollary}[冻结 escort 测度的正幂矩塌缩]\label{cor:xi-fixed-freezing-escort-positive-moment-collapse}
沿用定理 \ref{thm:xi-fixed-freezing-power-sum-ratio-locking} 的记号。对任意固定整数 \(k\ge 1\)，若
\[
X\sim \pi_{m,a},
\qquad
\pi_{m,a}(x):=\frac{d_m(x)^a}{S_a(m)},
\qquad
a>a_{\mathrm c},
\]
则有
\[
\boxed{\;
\EE_{\pi_{m,a}}\!\left[\left(\frac{d_m(X)}{M_m}\right)^k\right]
=
1+O\bigl(e^{-m\Delta(a)+o(m)}\bigr),\;}
\]
并且
\[
\boxed{\;
\EE_{\pi_{m,a}}\!\left[\left(1-\frac{d_m(X)}{M_m}\right)^k\right]
=
O\bigl(e^{-m\Delta(a)+o(m)}\bigr).\;}
\]
\end{corollary}
\begin{proof}
由定义直接得到
\[
\EE_{\pi_{m,a}}\!\left[\left(\frac{d_m(X)}{M_m}\right)^k\right]
=
\frac{S_{a+k}(m)}{M_m^kS_a(m)}.
\]
代入定理 \ref{thm:xi-fixed-freezing-power-sum-ratio-locking} 即得第一式。

对第二式，作二项展开
\[
\left(1-\frac{d_m(X)}{M_m}\right)^k
=
\sum_{j=0}^{k}(-1)^j\binom{k}{j}\left(\frac{d_m(X)}{M_m}\right)^j.
\]
取 \(\pi_{m,a}\)-期望后，\(j=0\) 项恒等于 \(1\)，而对每个 \(1\le j\le k\) 由第一式有
\[
\EE_{\pi_{m,a}}\!\left[\left(\frac{d_m(X)}{M_m}\right)^j\right]
=
1+O\bigl(e^{-m\Delta(a)+o(m)}\bigr).
\]
于是
\[
\EE_{\pi_{m,a}}\!\left[\left(1-\frac{d_m(X)}{M_m}\right)^k\right]
=
\sum_{j=0}^{k}(-1)^j\binom{k}{j}
+O\bigl(e^{-m\Delta(a)+o(m)}\bigr).
\]
再利用
\[
\sum_{j=0}^{k}(-1)^j\binom{k}{j}=0
\]
便得结论。证毕。
\end{proof}

\begin{theorem}[零温缩放极限的原子--地基精确分裂]\label{thm:xi-zero-temp-atomic-ground-exact-splitting}
沿用附录 \ref{app:real-input-40-zeta-u} 的记号。对真实输入 \(40\) 状态核，
\[
\Delta(z,u)=(1-uz^2)\,F(z,u).
\]
再令
\[
B_\infty=
\begin{pmatrix}
0&1&1&0\\
0&0&1&0\\
0&0&3&1\\
1&0&0&3
\end{pmatrix}.
\]
则对任意固定 \(w\in\CC\)，在零温缩放
\[
z=\sqrt{w/u}
\]
下有
\[
\boxed{\;
\lim_{u\to+\infty}\Delta\!\left(\sqrt{w/u},u\right)
=
(1-w)\det(I-wB_\infty).\;}
\]
等价地，零温窗口内的极限解析对象严格分裂为一个 primitive/Big-Witt 原子因子 \((1-w)\) 与一个四状态地基 SFT 的动力学行列式 \(\det(I-wB_\infty)\)。
\end{theorem}
\begin{proof}
由附录 \ref{app:real-input-40-zeta-u}，
\[
\begin{aligned}
F(z,u)
&=
1-z-6uz^2+2uz^3+(9u^2-u)z^4\\
&\qquad +(u-3u^2)z^5-(u^3+2u^2)z^6\\
&\qquad +(4u^3-3u^2)z^7+(u^3-u^4)z^8.
\end{aligned}
\]
把
\[
z=\sqrt{w/u}
\]
代入，得到
\[
\begin{aligned}
F\!\left(\sqrt{w/u},u\right)
&=
1-u^{-1/2}w^{1/2}-6w+2u^{-1/2}w^{3/2}\\
&\qquad +(9-u^{-1})w^2+(u^{-3/2}-3u^{-1/2})w^{5/2}\\
&\qquad -(1+2u^{-1})w^3+(4u^{-1/2}-3u^{-3/2})w^{7/2}\\
&\qquad +(u^{-1}-1)w^4.
\end{aligned}
\]
令 \(u\to+\infty\)，所有带负半整数次幂的项都消失，于是
\[
\lim_{u\to+\infty}F\!\left(\sqrt{w/u},u\right)
=
1-6w+9w^2-w^3-w^4.
\]
另一方面，由定理 \ref{thm:xi-zero-temperature-fourstate-artin-mazur-zeta-primitive-asymptotic}，
\[
1-6w+9w^2-w^3-w^4=\det(I-wB_\infty).
\]
再由
\[
1-u\left(\sqrt{w/u}\right)^2=1-w,
\]
便得到
\[
\lim_{u\to+\infty}\Delta\!\left(\sqrt{w/u},u\right)
=
(1-w)\det(I-wB_\infty).
\]
证毕。
\end{proof}

\begin{corollary}[零温主极点由地基层而非原子层支配]\label{cor:xi-zero-temp-leading-pole-ground-dominance}
沿用定理 \ref{thm:xi-zero-temp-atomic-ground-exact-splitting} 的记号。记
\[
y_\infty:=\rho(B_\infty)\approx 3.596154676009,
\qquad
b_\infty:=y_\infty^{-1}\approx 0.278074802141.
\]
则极限多项式
\[
(1-w)\det(I-wB_\infty)
\]
在正实轴上的最内层零点为
\[
\boxed{\;w=b_\infty,\;}
\]
而非原子零点 \(w=1\)。因此，零温缩放窗口中的主导自由能尺度完全由四状态地基层的 Perron 根 \(y_\infty\) 决定，而不由原子因子 \((1-w)\) 决定。
\end{corollary}
\begin{proof}
由定理 \ref{thm:xi-zero-temperature-fourstate-artin-mazur-zeta-primitive-asymptotic}，
\[
\det(I-wB_\infty)=1-6w+9w^2-w^3-w^4.
\]
其零点恰为 \(B_\infty\) 各特征值的倒数。由于 \(B_\infty\) 为非负不可约矩阵，其 Perron 根
\[
y_\infty=\rho(B_\infty)>1
\]
给出唯一最小正实倒数
\[
b_\infty=y_\infty^{-1}\in(0,1).
\]
另一方面，原子因子 \((1-w)\) 的正实零点为 \(w=1\)。由于
\[
b_\infty<1,
\]
在全部正实零点中更靠近原点的是 \(w=b_\infty\) 而非 \(w=1\)。故主导尺度来自地基层行列式 \(\det(I-wB_\infty)\)，而非原子因子。证毕。
\end{proof}

\begin{theorem}[零温地基四次数域与 Lee--Yang 三次数域的线性互不相交]\label{thm:xi-zero-temp-quartic-leyang-cubic-linear-disjoint}
记
\[
K_\infty:=\QQ(y_\infty),
\qquad
y_\infty=\rho(B_\infty),
\]
其中 \(y_\infty\) 满足
\[
y^4-6y^3+9y^2-y-1=0.
\]
再令
\[
K_{\mathrm{LY}}:=\QQ(u)=\QQ(b_{\mathrm{LY}}),
\]
其中 \(u\) 与 \(b_{\mathrm{LY}}\) 分别为文中 Lee--Yang 分支系数与振幅平方常数所生成的共同三次数域。则有
\[
\boxed{\;
[K_\infty:\QQ]=4,\qquad [K_{\mathrm{LY}}:\QQ]=3,\qquad
K_\infty\cap K_{\mathrm{LY}}=\QQ.\;}
\]
从而
\[
\boxed{\;
[K_\infty K_{\mathrm{LY}}:\QQ]=12,\;}
\]
并且 \(K_\infty\) 与 \(K_{\mathrm{LY}}\) 在 \(\QQ\) 上线性互不相交。
\end{theorem}
\begin{proof}
由定理 \ref{thm:xi-zero-temperature-fourstate-perron-parry-closed-form}，多项式
\[
y^4-6y^3+9y^2-y-1
\]
在 \(\QQ\) 上不可约，故
\[
[K_\infty:\QQ]=4.
\]
另一方面，由推论 \ref{cor:xi-terminal-zm-leyang-cubic-rational-observable-primitivity}，
\[
K_{\mathrm{LY}}=\QQ(u)=\QQ(b_{\mathrm{LY}})
\]
且
\[
[K_{\mathrm{LY}}:\QQ]=3.
\]

设
\[
E:=K_\infty\cap K_{\mathrm{LY}}.
\]
由塔式公式，
\[
[E:\QQ]\mid [K_\infty:\QQ]=4,
\qquad
[E:\QQ]\mid [K_{\mathrm{LY}}:\QQ]=3.
\]
因此 \([E:\QQ]\) 同时整除 \(4\) 与 \(3\)，只能等于 \(1\)。于是
\[
K_\infty\cap K_{\mathrm{LY}}=\QQ.
\]
再由塔式公式，
\[
[K_\infty K_{\mathrm{LY}}:\QQ]
=
\frac{[K_\infty:\QQ][K_{\mathrm{LY}}:\QQ]}{[K_\infty\cap K_{\mathrm{LY}}:\QQ]}
=
4\cdot 3
=
12.
\]
由于特征 \(0\) 下有限扩张皆可分，乘积次数达到乘积上界即等价于线性互不相交，故结论成立。证毕。
\end{proof}

\endinput
